% Bagian: first-order-logic
% Bab: introduction
% Subbagian: syntax

\documentclass[../../../include/open-logic-section]{subfiles}

\begin{document}

\olfileid[id]{fol}{int}{syn}

\olsection{Sintaksis}

Pertama-tama kita harus menetapkan secara tepat untai simbol mana yang
tergolong !!{sentence}s logika orde pertama. Kita akan melakukannya nanti;
untuk sekarang, kita cukup melanjutkan melalui contoh. Unsur-unsur dasar---
kosakata---logika orde pertama terbagi menjadi dua bagian. Bagian pertama
terdiri atas simbol-simbol yang kita gunakan untuk mengatakan sesuatu yang
tertentu atau untuk menunjuk sesuatu yang tertentu. Kita menunjuk sesuatu
dengan menggunakan !!{constant}s, dan kita mengatakan sesuatu mengenai objek
yang kita tunjuk dengan menggunakan !!{predicate}s. Misalnya, kita dapat
menggunakan $\Obj a$ sebagai !!a{constant} untuk menunjuk satu objek, lalu
mengatakan sesuatu mengenainya dengan menggunakan !!{sentence}~$\Atom{\Obj
P}{\Obj a}$. Jika Anda membayangkan makna bagi ``$\Obj a$'' dan ``$\Obj P$'',
Anda dapat membaca $\Atom{\Obj P}{\Obj a}$ sebagai suatu kalimat bahasa Inggris
(dan Anda mungkin pernah melakukannya ketika pertama kali mempelajari logika
formal). Setelah memiliki !!{sentence}s logika orde pertama yang sederhana
semacam itu, Anda dapat membangun kalimat-kalimat yang lebih kompleks dengan
menggunakan bagian kedua kosakata tersebut: simbol logis (penghubung dan
kuantor). Jadi, misalnya, kita dapat membentuk ungkapan seperti
$(\Atom{\Obj P}{\Obj a} \land \Atom{\Obj Q}{\Obj b})$ atau~$\lexists[\Obj x][\Atom{\Obj P}{\Obj x}]$.

Untuk memberikan definisi semantik yang tepat dan aturan-aturan sistem
!!{derivation} kita yang diperlukan bagi kajian metalogis yang ketat,
pertama-tama kita harus memberikan definisi yang tepat tentang apa yang
tergolong !!a{sentence} logika orde pertama. Gagasan dasarnya cukup mudah
dipahami: ada beberapa !!{sentence}s sederhana yang dapat kita bentuk hanya
dari !!{predicate}s dan !!{constant}s, seperti~$\Atom{\Obj P}{\Obj a}$. Lalu,
dari kalimat-kalimat ini kita membentuk kalimat yang lebih kompleks dengan
menggunakan penghubung dan kuantor. Namun, persis aturan apa yang mengizinkan
kita membentuk !!{sentence}s yang lebih kompleks? Aturan-aturan itu harus
ditetapkan; jika tidak, kita belum mendefinisikan ``!!{sentence} logika orde
pertama'' dengan cukup tepat. Ada beberapa persoalan. Persoalan pertama ialah
menentukan untai yang tepat agar tergolong !!{sentence}s. Persoalan kedua
ialah melakukannya sedemikian rupa sehingga kita dapat memberikan bukti
matematis mengenai \emph{semua} !!{sentence}s. Terakhir, kita juga harus
memberikan definisi yang tepat bagi sejumlah operasi dasar atas !!{sentence}s,
seperti ``ganti setiap $\Obj x$ dalam~$!A$ dengan~$\Obj b$.'' Kesulitannya,
kuantor dan !!{variable}s dalam logika orde pertama membuat cara melakukannya
tidak sepenuhnya jelas. Misalnya, apakah $\lexists[\Obj x][\Atom{\Obj P}{\Obj
a}]$ tergolong !!a{sentence}? Bagaimana dengan $\lexists[\Obj x][\lexists[\Obj
x][\Atom{\Obj P}{\Obj x}]]$? Apa hasil dari ``ganti $\Obj x$ dengan~$\Obj b$
dalam $(\Atom{\Obj P}{\Obj x} \land \lexists[\Obj x][\Atom{\Obj P}{\Obj
x}])$''?

\end{document}
